\section{Pipeline de Dados e Artefatos}

O pipeline Python é responsável por transformar os dados externos em dados confiáveis para a aplicação. O processamento utiliza DataFrames e preserva arquivos Parquet como artefatos intermediários.

\subsection{Camadas de dados}

A organização conceitual é:

\begin{lstlisting}
SINAN / PySUS
    |
    v
  Bronze
    |
    v
  Silver
    |
    v
   Gold
    |
    +----> PostgreSQL
\end{lstlisting}

\begin{itemize}
    \item \textbf{Bronze}: preserva os dados de entrada e informações da ingestão;
    \item \textbf{Silver}: concentra limpeza, padronização e transformação;
    \item \textbf{Gold}: contém dados consolidados e agregações adequadas às consultas do sistema.
\end{itemize}

\subsection{Papel do DataFrame}

O DataFrame é uma estrutura de processamento do pipeline. Ele facilita transformações, validações, agrupamentos e inspeção durante a execução.

\subsection{Papel do Parquet}

Os arquivos Parquet são mantidos como artefatos persistentes do pipeline. Eles são úteis para:

\begin{itemize}
    \item inspeção manual dos resultados;
    \item comparação entre etapas;
    \item reprodução de uma transformação sem coletar novamente os dados externos;
    \item investigação de falhas;
    \item auditoria técnica da evolução dos dados.
\end{itemize}

\begin{explicacao}[Parquet não substitui o PostgreSQL]
O Parquet continua sendo parte importante do pipeline, mas não é o armazenamento principal utilizado pelo frontend. O PostgreSQL funciona como fonte persistida dos dados consolidados para a aplicação Java.
\end{explicacao}

\subsection{API de inspeção do pipeline}

O componente Python pode disponibilizar uma API de inspeção para desenvolvimento, validação e debugging. Essa API não deve ser confundida com a API oficial da aplicação.

Seu objetivo é permitir comparar diretamente o resultado do pipeline com o que posteriormente será armazenado no PostgreSQL e retornado pelo Java.

\subsection{Validação antes da publicação}

Antes de considerar uma execução concluída, o pipeline deve validar propriedades importantes dos dados. Exemplos:

\begin{itemize}
    \item número de registros por etapa;
    \item presença das colunas esperadas;
    \item valores obrigatórios;
    \item totais por doença, ano e unidade geográfica;
    \item consistência entre Gold e dados persistidos no PostgreSQL.
\end{itemize}

Uma falha de validação deve ser registrada e impedir a publicação de dados inconsistentes, quando a regra de validação for crítica.
